\newpage
\section{Limitations and next steps}
The optimized choices apply to the declared family and tested workload. The controlled
distribution creates a clear score gap: twelve coordinates dominate all others combined.
That is a useful test of the indexing idea, not evidence that typical text embeddings have
the same structure. Matryoshka prefixes alone do not establish the assumption.

The cost model is more reliable for the selected operating points than across every setting.
Memory access order, result-selection behavior and different query mixtures can change
latency even when aggregate work counts match. The new occupancy terms improve the estimate
without making it a universal hardware model. The next relevant check on real embeddings
is whether a small, query-dependent coordinate set actually determines their nearest neighbors.

A larger collection does not automatically create a crossover. Under the simple fixed-depth
model, write $T_A=t_A+fNc_s$ and $T_B=t_B+N\beta_B$, where $f$ is A's scanned fraction and
$\beta_B=(mc_b+c_l)/w+2^{-m}c_g$. Their candidate crossing is
\[
 N^*=\frac{t_B-t_A}{fc_s-\beta_B}.
\]
It is useful only within the model's valid range and while recall and RAM constraints hold.
If B has both a higher fixed cost and a larger per-document cost, increasing $N$ creates no
win in this model. At one billion 256-bit documents, codes and eight-byte IDs alone need 40 GB;
B adds at least 32 GB of bitplanes. All current layouts therefore exceed a 32 GB budget at
that size, before program and temporary memory. No billion-document experiment is claimed.

\section{Conclusion}
Scanning was fastest in the real-data replication. In the constructed fixed-position case,
backward walk was about 10\% faster than a direct-prefix scan that selected the same candidate
group. With changing important coordinates, bitplanes were about 152 times faster than the
best tested scan. These results support different choices for different workloads; they do
not establish a universal winner or the prevalence of those conditions in real embeddings.

The formulas identify candidates and explain their work. The final decision still combines
measured recall, memory and complete retrieval time, with every method allowed its own
settings. The selected-point prediction check, the retained failures and the real-data
replication define the scope of that conclusion.

\begin{thebibliography}{9}\small\raggedright
\bibitem{exa} Exa Team. \emph{How we built a web-scale vector database}. 2024.
\url{https://exa.ai/blog/building-web-scale-vector-db}.
\bibitem{mrl} A. Kusupati et al. \emph{Matryoshka Representation Learning}. 2022.
\url{https://arxiv.org/abs/2205.13147}.
\bibitem{msmarco} Microsoft. \emph{MS MARCO}. Passage collection and benchmark resources.
\url{https://microsoft.github.io/msmarco/}.
\bibitem{nomic} Nomic AI. \emph{nomic-embed-text-v1.5}. Model documentation.
\url{https://huggingface.co/nomic-ai/nomic-embed-text-v1.5}.
\end{thebibliography}
